\subsection{Alternative investment structures, fees and returns}

\begin{qfdefinition}[Alternative investments: features and categories, and direct investment, co-investment and fund investment compared]\label{def:alternative-investment-categories}
\qffield{Statement} \emph{Categories.} ``Alternative'' covers investments other than long-only public stocks, bonds and cash. The main families are (Interpretation, organised from the sources below):
\begin{itemize}
\item \emph{private capital}: private equity (venture, growth, buyout, turnaround) and private debt (direct lending, mezzanine, distressed debt) \citep[pp.~58--60]{demaria-private-markets-2021};
\item \emph{real assets}: real estate, infrastructure and natural resources such as timberland and farmland, held directly or through funds and listed vehicles (pp.~58--59);
\item \emph{hedge funds}, which trade mostly listed securities and derivatives under flexible mandates \citep[pp.~559--562]{s7-ang};
\item \emph{commodities}, usually through futures, and \emph{digital assets}.
\end{itemize}

\emph{Common features.} Relative to public securities, the private families are illiquid, traded over the counter with large transaction costs, valued infrequently and subjectively, held for about ten years under complex contracts \citep[p.~594, Table~18.1]{s7-ang}. Investors depend on specialist managers whose pay includes performance fees. Return data are voluntary and biased by survivorship and selection, and reported returns are smoothed (pp.~416--426). Private-equity investors also take board seats and control rights to influence management, which hedge funds and mutual funds do not, and they need an exit to return cash \citep[pp.~5--7]{metrick-yasuda-survey-2011}.

\emph{Methods of investing} \citep[pp.~60--66, 95--98]{demaria-private-markets-2021}:
\begin{itemize}
\item \emph{Fund investment.} The investor commits to a pooled vehicle run by a manager (GP) who selects, monitors and exits the assets. It gives diversification and expertise for modest resources, at the price of fees, carried interest and little control. Funds of funds add a further layer of diversification and fees.
\item \emph{Co-investment.} The investor invests directly in a deal alongside a fund, often at reduced or no fees, and can overweight particular assets. It needs staff able to decide within weeks. It risks adverse selection if the manager offers only deals that exceed the fund's own capacity. The return correlation of LBO co-investments with the funds rises with deal size, from 0.42 for deals up to \$25 million to 0.87 above \$2.5 billion (Table~3.1, 2006--2016).
\item \emph{Direct investment.} The investor sources, owns and manages assets itself. It removes fees and gives full control but requires teams that can compete with fund managers, enough capital to diversify, and freedom from home bias. Institutions with seven large direct programmes outperformed only where they could resolve information problems (p.~67, citing Fang, Ivashina and Lerner).
\end{itemize}

\qffield{Practice} \emph{Choosing a method} (original). A \$2 billion foundation with a two-person private-markets team should use funds, perhaps a fund of funds for venture, because it cannot diversify or staff direct deals. A \$60 billion pension with an in-house team can add systematic co-investment rights to its large buyout commitments to cut its average fee, provided it accepts the concentration. Direct ownership suits assets it understands and can hold for decades, such as core infrastructure.

\qffield{Conditions} Sizes and correlations are as reported by the sources for their samples; the World Economic Forum thresholds Demaria cites (co-investing above about \$25 billion of assets, direct investing above \$50 billion) are rules of thumb, not requirements.

\qffield{Pitfalls} Treating co-investment as free alpha without the resources to select deals. Comparing the fee saving of direct investment without the cost of the team. Counting a hedge fund as private capital because both charge performance fees.

\qffield{See also} \hyperref[def:alternative-fund-terms]{fund terms and fees}; \hyperref[meth:private-fund-structures]{fund structures}; \hyperref[def:alternatives-roles-opportunity-set]{roles in a portfolio}.

\end{qfdefinition}

\begin{qfdefinition}[Ownership and compensation structures in alternative investments: limited partnerships, management fees, carried interest and incentive fees, hurdles, catch-up, high-water marks, clawbacks and waterfalls]\label{def:alternative-fund-terms}
\qffield{Statement} \emph{Ownership.} Private funds are usually limited partnerships with a fixed term of about ten years. Investors are \emph{limited partners} (LPs), whose liability is limited to their commitment; the manager is the \emph{general partner} (GP), with management usually delegated to a management company that receives the fee \citep[pp.~13--15]{ippolito-private-capital-2020}. LPs commit capital that is called as deals arise; once committed they are expected to stay to the end \citep[p.~594]{s7-ang}. Hedge funds are open-ended, with periodic subscriptions and redemptions and no fixed life \citep[pp.~6--7]{metrick-yasuda-survey-2011}.

\emph{Management fees.}
\begin{itemize}
\item Private funds: typically 1.0 percent (private debt) to 2.0 percent (private equity) a year of \emph{committed} capital during the investment period, then on the cost of unrealised investments, often stepped down by 0.5--1.0 points \citep[pp.~15--16]{ippolito-private-capital-2020}. Charging on commitments discourages a rush to invest but pays the GP on money not yet working \citep[p.~611]{s7-ang}.
\item Hedge funds: typically 2 percent a year of net asset value \citep[p.~581]{s7-ang}.
\end{itemize}

\emph{Performance pay.}
\begin{itemize}
\item \emph{Carried interest} in private funds, typically 20 percent of profits, is earned only on realisations, after LPs have their capital back and a \emph{hurdle} or \emph{preferred return}, customarily 7--8 percent a year for private equity and 3--4 percent for private debt \citep[p.~16]{ippolito-private-capital-2020}.
\item A \emph{catch-up} then gives the GP all or most distributions until it has its 20 percent of total profit; thereafter distributions split 80/20 (p.~16; Ang, p.~611).
\item A \emph{clawback} returns to LPs any carry above the agreed share of cumulative profit (Ang, p.~611).
\item The order of payments is the \emph{waterfall}. A whole-of-fund waterfall, called European in market usage, pays carry only when the fund as a whole has returned capital and hurdle. A deal-by-deal (American) waterfall pays carry on each exit and relies on the clawback when later deals lose (Interpretation). Ippolito notes that deal-by-deal carry ``has nearly disappeared'' \citep[p.~16]{ippolito-private-capital-2020}.
\item \emph{Hedge-fund incentive fees}, typically 20 percent, are charged each period on gains, often above a hurdle and subject to a \emph{high-water mark}: no fee until net asset value exceeds its previous peak \citep[pp.~581--582]{s7-ang}. Unlike private-equity carry, this is charged on unrealised marks.
\item \emph{Other fees}: transaction, monitoring and advisory fees charged to portfolio companies, best offset against management fees \citep[pp.~611--612]{s7-ang}.
\end{itemize}

\emph{Economics.} The incentive fee is a call option on fund performance and encourages risk-taking \citep[pp.~584--585]{s7-ang}. Across 238 funds raised in 1993--2006, about two-thirds of expected GP revenue was fixed (fees) rather than performance-linked \citep[p.~35]{metrick-yasuda-survey-2011}.

\qffield{Practice} \emph{Ashby Fund I} (original; \hyperref[meth:gp-alignment-net-returns]{numbers}). A buyout fund with a 2 percent fee on 100 committed, an 8 percent compounding hurdle, full catch-up and 20 percent carry pays its GP 11.85 of carry, exactly 20 percent of the fund's profit of 59.25, all in year 7. On a deal-by-deal waterfall the GP would have received 13.21, starting in year 4, and would owe 1.36 back under the clawback.

\qffield{Conditions} Terms are negotiated fund by fund and recorded in the limited partnership agreement or offering documents; the figures above are typical ranges reported by the sources, not standards.

\qffield{Pitfalls} Assuming a 20 percent carry means a 20 percent share of profit above the hurdle: with a full catch-up it is 20 percent of all profit. Ignoring the creditworthiness of a GP when relying on a clawback. Comparing hedge-fund and private-equity ``2 and 20'' as if charged on the same base.

\qffield{See also} \hyperref[meth:alternative-returns-after-fees]{returns after fees}; \hyperref[meth:gp-alignment-net-returns]{GP alignment}; \hyperref[def:hedge-fund-features]{hedge funds}.

\end{qfdefinition}

\begin{qfmethod}[Appraising alternative investment performance and calculating returns before and after fees: return measures, data biases, smoothing, and hedge-fund and private-fund fee calculations]\label{meth:alternative-returns-after-fees}
\qffield{Statement} \emph{Appraisal.} Liquid strategies are measured by time-weighted returns and risk-adjusted measures (\hyperref[def:return-measures]{return measures}). Private funds are measured by dated cash flows: the internal rate of return, multiples of invested capital and public-market equivalents \citep[pp.~597--602]{s7-ang}. Four problems distort the reported figures (pp.~416--426):
\begin{itemize}
\item \emph{survivorship and reporting bias}: funds that fail stop reporting, and funds with poor records never start, so averages are too high;
\item \emph{infrequent trading and appraisal}: values move with a lag, so volatility, correlation and beta are understated and Sharpe ratios overstated;
\item \emph{selection bias}: assets are sold, and so priced, when values are high;
\item \emph{IRR manipulation}: early distributions raise the IRR without raising wealth (pp.~598--600).
\end{itemize}
Returns are compared net of all fees, the only measure relevant to the investor. Gross returns describe the manager's value added \citep[p.~67]{demaria-private-markets-2021}.

\emph{Calculation.} For a hedge fund over each period:
\[ V^{\mathrm{after\ mgmt}}=V_0(1+r_{\mathrm{gross}})-m\,V_0,\qquad \text{incentive}=p\max\{0,\;V^{\mathrm{after\ mgmt}}-H(1+h)\}, \]
where $m$ is the management fee, $p$ the incentive rate, $H$ the high-water mark and $h$ the hurdle; the net return is the closing value over $V_0$, less one (Derivation, from the fee terms above). Conventions differ on whether fees accrue monthly, whether the hurdle compounds, and whether a \emph{soft} hurdle, once cleared, charges the fee on the whole gain; the offering documents decide. For a private fund, carry follows the waterfall and the net IRR is computed from the LP's dated calls (including fees) and distributions.

\qffield{Practice} \emph{Brook Relative Value} (original). Gross returns 15, $-10$, 12 and 20 percent; fees 2/20 with a high-water mark and a 4 percent hurdle on it; fees charged at year end on opening NAV:
\begin{center}\small
\begin{tabular}{@{}lrrrr@{}}
\toprule
Year & 1 & 2 & 3 & 4 \\
\midrule
Gross return & 15.00\% & $-10.00\%$ & 12.00\% & 20.00\% \\
Management fee & 2.00 & 2.22 & 1.96 & 2.15 \\
Incentive fee & 1.80 & 0.00 & 0.00 & 2.27 \\
Closing NAV & 111.20 & 97.86 & 107.64 & 124.74 \\
Net return & 11.20\% & $-12.00\%$ & 10.00\% & 15.89\% \\
\bottomrule
\end{tabular}
\end{center}
\begin{itemize}
\item In year 1 the fee applies to $113.00-104.00=9.00$. In year 3 the fund is below its mark of 111.20 and pays none; in year 4 it pays on $127.02-115.65$.
\item Over four years 100 grows to 139.10 gross and 124.74 net: 8.60 against 5.68 percent a year, a fee drag of 2.92 points.
\item Without the high-water mark year 3 would pay 1.17 for recovering earlier losses. With a soft hurdle year 1 would pay 2.60.
\end{itemize}
\emph{Ashby Fund I} (\hyperref[meth:gp-alignment-net-returns]{details}): deals earn a gross IRR of 16.87 percent and 1.88 times cost; LPs receive 11.71 percent and 1.52 times paid-in after fees of 10.75 and carry of 11.85.

\qffield{Conditions} Annualising a since-inception IRR, or comparing it with a time-weighted index return, mixes measures with different weights on time. Smoothing corrections assume a smoothing model (\hyperref[meth:private-capital-diversification]{diversification}).

\qffield{Pitfalls} Quoting gross returns against net benchmarks. Treating an index of surviving funds as investable. Reading a high Sharpe ratio from appraisal-based returns as low risk.

\qffield{See also} \hyperref[def:alternative-fund-terms]{fund terms}; \hyperref[meth:private-fund-multiples-pme]{fund multiples and PME}; \hyperref[def:return-measures]{return measures}.

\end{qfmethod}
